\section{Processi Primari}
\subsubsection{Strumenti a supporto}
Il gruppo si avvale dei seguenti strumenti:
\begin{itemize}
	\item \textbf{Email del gruppo}: l'indirizzo {\emailgruppo} viene utilizzato per tutte le comunicazioni ufficiali con committente e proponente.
	\item \textbf{Repository GitHub}: punto di raccolta, versionamento e pubblicazione di tutta la documentazione e del codice del POC;
	\item \textbf{GitHub Pages}: utilizzato per esporre la documentazione prodotta;
\end{itemize}

\subsubsection{Documentazione prodotta in fase di Candidatura}
Durante la fase di candidatura, il gruppo si è occupato di produrre e approvare tutta la documentazione necessaria per la presentazione della propria proposta.

\begin{table}[htbp]
	\centering
	\caption{Registro documentazione prodotta}
	\label{tab:documenti}
	\renewcommand{\arraystretch}{1.3}
	\small
	\begin{tabularx}{\textwidth}{|>{\centering\arraybackslash}p{4.5cm}
			|X|}
		\hline
		\rowcolor[gray]{0.88}
		\textbf{Documento} &
		\centering\arraybackslash\textbf{Descrizione} \\
		\hline
		Lettera di Presentazione & Documento con cui il gruppo comunica in maniera formale la candidatura al capitolato selezionato. Contiene le informazioni sul capitolato, sulla proponente, sulla scadenza fissata, del budget preventivato e gli impegni assunti. \\ \hline
		Valutazione dei Capitolati & Documento che contiene il confronto tra tutti i capitolati presentati, analizzati da parte del gruppo secondo aspetti positivi, criticità e rischi. \\ \hline
		Preventivo Costi e Assunzione Impegni & Documento che formalizza il preventivo economico, fissa le scadenze, la distribuzione delle ore produttive e la ripartizione dei ruoli.  \\ \hline
		Norme di Progetto & Documento interno contenente le regole operative. \\ \hline
		Verbali interni & Documento interno il cui scopo è tracciare le decisioni organizzative e le assegnazioni delle attività del gruppo. \\ \hline
		Verbali esterni & Documento esterno che formalizza gli incontri con le aziende, gli argomenti discussi e i chiarimenti ottenuti. \\ \hline
	\end{tabularx}
\end{table}

\subsection{Incontri con la proponente}
Per la richiesta di incontri con la proponente il gruppo ha deciso di procedere nel seguente modo:
\begin{enumerate}
	\item Il gruppo richiede l'incontro per email provvedendo a inviare in anticipo l'ordine del giorno ed eventuali domande;
	\item Prima dell'incontro il gruppo dovrà già aver assegnato il ruolo di scriba;
	\item Al termine dell'incontro lo scriba si occuperà della stesura del verbale esterno, che dovrà essere verificato da un Verificatore e infine viene approvato e firmato dal responsabile per il rilascio;
	\item Quando il documento raggiunge la versione 1.0.0, il verbale verrà inviato al proponente che ne firma l'ultima sezione
	\item Il verbale dopo essere stato firmato dalla proponente verrà caricato nella repository nella cartella \texttt{Verbali esterni}.
\end{enumerate}